%% ╔══════════════════════════════════════════════════════════════════╗
%% ║         CORIOTLAB — PRESENTACIÓN ACADÉMICA / PROFESIONAL        ║
%% ║         LaTeX Beamer · diseño basado en Brand Book CORIOTLAB    ║
%% ╚══════════════════════════════════════════════════════════════════╝

%% ═══════════════════════════════════════════════════════════════════
%% ZONA DE CONFIGURACIÓN
%% ═══════════════════════════════════════════════════════════════════
\newcommand{\PresTitle}{Sistema de Monitoreo IoT Distribuido}
\newcommand{\PresSubtitle}{Red de sensores ESP32 con MQTT, InfluxDB y Grafana}
\newcommand{\PresAutor}{Santiago Leal}
\newcommand{\PresCargo}{Investigador --- CORIOTLAB}
\newcommand{\PresProyecto}{Proyecto de Monitoreo Ambiental}
\newcommand{\PresFecha}{19 de junio de 2025}
\newcommand{\PresInstitucion}{CORIOTLAB --- ITM}
%% ═══════════════════════════════════════════════════════════════════

\documentclass[aspectratio=169, 12pt]{beamer}

%% --- Lenguaje
\usepackage[spanish, es-nodecimaldot, es-noshorthands]{babel}

%% --- Fuentes
\usepackage{fontspec}
\setmainfont{Inter}
\newfontfamily\FuenteTitulo{MuseoModerno}
\newfontfamily\FuenteCodigo{Space Mono}

%% --- Color
\usepackage{xcolor}
\definecolor{AzulITM}     {HTML}{102D69}
\definecolor{Magenta}     {HTML}{C14894}
\definecolor{AzulDigital} {HTML}{56ACDE}
\definecolor{GrisPizarra} {HTML}{2F2F2F}
\definecolor{GrisClaro}   {HTML}{F2F2F2}
\definecolor{GrisMedio}   {HTML}{AAAAAA}
\definecolor{GrisLinea}   {HTML}{DDDDDD}
\definecolor{Blanco}      {HTML}{FFFFFF}

%% --- Gráficos y TikZ
\usepackage{tikz}
\usepackage{graphicx}

%% --- Cajas
\usepackage[most]{tcolorbox}

%% --- Código
\usepackage{listings}

%% --- Tablas
\usepackage{array}
\usepackage{colortbl}

%% ─────────────────────────────────────────────────────────────────
%% TEMA BEAMER
%% ─────────────────────────────────────────────────────────────────
\usetheme{default}
\useinnertheme{default}
\useoutertheme{default}

\setbeamercolor{normal text}     {fg=GrisPizarra, bg=white}
\setbeamercolor{structure}       {fg=AzulITM}
\setbeamercolor{palette primary} {bg=AzulITM, fg=white}
\setbeamercolor{frametitle}      {bg=AzulITM, fg=white}
\setbeamercolor{title}           {fg=white}
\setbeamercolor{subtitle}        {fg=AzulDigital}
\setbeamercolor{block title}     {bg=AzulITM, fg=white}
\setbeamercolor{block body}      {bg=GrisClaro, fg=GrisPizarra}
\setbeamercolor{alerted text}    {fg=Magenta}
\setbeamercolor{item}            {fg=AzulITM}
\setbeamercolor{footline}        {fg=GrisMedio, bg=white}
\setbeamercolor{enumerate item}  {fg=AzulITM}

\setbeamerfont{title}      {family=\FuenteTitulo, size=\huge,      series=\bfseries}
\setbeamerfont{subtitle}   {family=\FuenteTitulo, size=\large,     series=\mdseries}
\setbeamerfont{frametitle} {family=\FuenteTitulo, size=\large,     series=\bfseries}
\setbeamerfont{block title}{family=\FuenteTitulo, size=\normalsize, series=\bfseries}
\setbeamerfont{footline}   {size=\scriptsize}

%% Bullets Beamer nativos con >
\setbeamertemplate{itemize item}   {\color{AzulITM}\textbf{>}}
\setbeamertemplate{itemize subitem}{\color{AzulDigital}\small\textbf{>}}
\setbeamertemplate{enumerate item} {\color{AzulITM}\textbf{\insertenumlabel.}}
\setbeamertemplate{navigation symbols}{}

\setbeamertemplate{frametitle}{%
  \vspace{0.3em}\insertframetitle\vspace{0.2em}}

%% ─────────────────────────────────────────────────────────────────
%% FOOTLINE
%% ─────────────────────────────────────────────────────────────────
\setbeamertemplate{footline}{%
  \leavevmode%
  \begin{beamercolorbox}[wd=\paperwidth, ht=0.6pt, dp=0pt]{structure}%
  \end{beamercolorbox}%
  \begin{beamercolorbox}[wd=\paperwidth, ht=2.0em, dp=0.8em,
      leftskip=1em, rightskip=1em]{footline}%
    {\FuenteTitulo\tiny\color{AzulITM}\textbf{< CORIOTLAB}}%
    \hfill%
    {\tiny\color{GrisMedio}\PresTitle}%
    \hfill%
    {\tiny\color{GrisMedio}\insertframenumber\,/\,\inserttotalframenumber}%
  \end{beamercolorbox}%
}
\setbeamertemplate{headline}{}

%% ─────────────────────────────────────────────────────────────────
%% CÓDIGO lstlisting
%% ─────────────────────────────────────────────────────────────────
\lstdefinestyle{coriotcode}{
  basicstyle=\FuenteCodigo\scriptsize,
  backgroundcolor=\color{GrisClaro},
  commentstyle=\color{GrisMedio}\itshape,
  keywordstyle=\color{AzulITM}\bfseries,
  stringstyle=\color{Magenta},
  numberstyle=\tiny\color{GrisMedio},
  numbers=left, numbersep=8pt,
  breaklines=true, frame=none,
  tabsize=2, showstringspaces=false
}
\lstset{style=coriotcode}

%% ─────────────────────────────────────────────────────────────────
%% CAJAS tcolorbox
%% ─────────────────────────────────────────────────────────────────
\tcbset{
  cajaNota/.style={
    colback=AzulDigital!10, colframe=AzulDigital,
    fonttitle=\FuenteTitulo\small\bfseries, coltitle=white,
    boxrule=0pt, leftrule=4pt, title=Nota,
    left=6pt, top=4pt, bottom=4pt
  },
  cajaCode/.style={
    colback=GrisClaro, colframe=AzulDigital,
    boxrule=0pt, leftrule=4pt,
    left=6pt, top=4pt, bottom=4pt
  },
  cajaIndicador/.style={
    colback=GrisClaro, colframe=GrisLinea,
    boxrule=0.5pt, halign=center,
    left=4pt, right=4pt, top=6pt, bottom=6pt
  }
}

%% ─────────────────────────────────────────────────────────────────
%% MACRO: LOGO CORIOTLAB en TikZ
%% ─────────────────────────────────────────────────────────────────
\newcommand{\LogoCoriotTikZ}[2]{%
  \begin{tikzpicture}[baseline=-0.3ex]
    \node[anchor=base east,  text=#1, font=\FuenteTitulo\Large\bfseries]
      at (0,0)     {<};
    \node[anchor=base west,  text=#2, font=\FuenteTitulo\Large\bfseries]
      at (0.05cm,0){CORIOTLAB};
    \fill[AzulDigital] (-0.85cm, 0.72cm) rectangle (-0.65cm, 0.52cm);
    \fill[Magenta]     (-0.65cm, 0.48cm) rectangle (-0.48cm, 0.32cm);
    \fill[AzulITM]     (-0.85cm, 0.28cm) rectangle (-0.72cm, 0.16cm);
  \end{tikzpicture}%
}

%% =================================================================
\begin{document}

%% -----------------------------------------------------------------
%% DIAPOSITIVA 1: PORTADA  [plain → sin footline]
%% -----------------------------------------------------------------
\begin{frame}[plain]
  \begin{tikzpicture}[remember picture, overlay]
    \fill[AzulITM] (current page.south west) rectangle (current page.north east);
    \fill[AzulDigital]
      ([xshift=-2.4cm, yshift=-1.0cm] current page.north east)
      rectangle ([xshift=-1.9cm, yshift=-1.5cm] current page.north east);
    \fill[Magenta]
      ([xshift=-1.8cm, yshift=-1.5cm] current page.north east)
      rectangle ([xshift=-1.4cm, yshift=-1.9cm] current page.north east);
    \fill[AzulDigital!60]
      ([xshift=-2.4cm, yshift=-2.0cm] current page.north east)
      rectangle ([xshift=-2.0cm, yshift=-2.4cm] current page.north east);
    \fill[Magenta!70]
      ([xshift=-1.3cm, yshift=-1.0cm] current page.north east)
      rectangle ([xshift=-1.0cm, yshift=-1.3cm] current page.north east);
    \fill[AzulDigital]
      (current page.south west) rectangle ([yshift=0.5cm] current page.south east);
  \end{tikzpicture}

  \vspace{1.2cm}
  \begin{center}
    \LogoCoriotTikZ{AzulDigital}{white}\\[1.2em]
    {\FuenteTitulo\huge\color{white}\textbf{\PresTitle}}\\[0.5em]
    {\color{AzulDigital}\rule{0.6\textwidth}{1.5pt}}\\[0.5em]
    {\FuenteTitulo\large\color{AzulDigital}\PresSubtitle}\\[1.4em]
    {\small\color{white}\PresAutor\ \ \textbullet\ \ \PresCargo}\\[0.2em]
    {\small\color{white}\PresInstitucion\ \ \textbullet\ \ \PresFecha}
  \end{center}
\end{frame}

%% -----------------------------------------------------------------
%% DIAPOSITIVA 2: AGENDA
%% -----------------------------------------------------------------
\begin{frame}{Agenda}
  \begin{tikzpicture}[overlay, remember picture]
    \fill[AzulITM]
      ([xshift=0.6cm,  yshift=-0.85cm] current page.north west) rectangle
      ([xshift=0.86cm, yshift=1.1cm]   current page.south west);
  \end{tikzpicture}

  \hspace{0.7cm}
  \begin{minipage}{0.88\textwidth}
    \begin{enumerate}
      \item Introducción y contexto del proyecto
      \item Arquitectura del sistema IoT
      \item Implementación del firmware ESP32
      \item Resultados y métricas de desempeño
      \item Conclusiones y trabajo futuro
    \end{enumerate}
  \end{minipage}
\end{frame}

%% -----------------------------------------------------------------
%% DIAPOSITIVA 3: SEPARADOR DE SECCIÓN  [plain]
%% -----------------------------------------------------------------
\begin{frame}[plain]
  \begin{tikzpicture}[remember picture, overlay]
    \fill[AzulITM] (current page.south west) rectangle (current page.north east);
    %% Símbolo decorativo "<" grande — opacidad 15% con scalebox
    \node[opacity=0.15, anchor=center] at (current page.center)
      {\scalebox{16}{\color{white}\FuenteTitulo<}};
    \node[anchor=south west] at ([xshift=2.2cm, yshift=1.2cm] current page.south west)
      {\FuenteTitulo\Huge\color{AzulDigital}\textbf{01}};
    \node[anchor=south west] at ([xshift=2.2cm, yshift=3.2cm] current page.south west)
      {\FuenteTitulo\Large\color{white}\textbf{Introducción y Contexto}};
  \end{tikzpicture}
\end{frame}

%% -----------------------------------------------------------------
%% DIAPOSITIVA 4: CONTENIDO ESTÁNDAR
%% -----------------------------------------------------------------
\begin{frame}{Arquitectura del Sistema de Monitoreo IoT}
  \begin{tikzpicture}[overlay, remember picture]
    \fill[AzulITM]
      ([xshift=0.6cm,  yshift=-0.85cm] current page.north west) rectangle
      ([xshift=0.85cm, yshift=1.1cm]   current page.south west);
  \end{tikzpicture}

  \hspace{0.55cm}
  \begin{minipage}{0.88\textwidth}
    La arquitectura sigue un patrón de tres capas:
    \vspace{0.4em}
    \begin{itemize}
      \item \textbf{Capa de percepción:} 8 nodos ESP32 con sensores BME280
        (temperatura, humedad, presión).
      \item \textbf{Capa de transporte:} Protocolo MQTT sobre WiFi 2.4\,GHz
        con broker HiveMQ CE local.
      \item \textbf{Capa de aplicación:} InfluxDB para almacenamiento
        + Grafana para visualización en tiempo real.
    \end{itemize}

    \vspace{0.5em}
    \begin{tcolorbox}[cajaNota, top=3pt, bottom=3pt]
      Cada nodo publica datos cada 5\,s con QoS 1, garantizando entrega
      incluso ante reconexiones de red.
    \end{tcolorbox}
  \end{minipage}
\end{frame}

%% -----------------------------------------------------------------
%% DIAPOSITIVA 5: DOS COLUMNAS
%% -----------------------------------------------------------------
\begin{frame}{Diagrama de Nodos de Sensado}
  \begin{columns}[T]
    \begin{column}{0.52\textwidth}
      \textbf{\color{AzulITM}Distribución en el laboratorio:}
      \begin{itemize}
        \item 3 nodos --- Zona de trabajo (banco norte)
        \item 2 nodos --- Zona de servidores
        \item 2 nodos --- Sala de reuniones
        \item 1 nodo --- Entrada / control de acceso
      \end{itemize}
      \vspace{0.4em}
      \textbf{\color{AzulITM}Parámetros monitoreados:}
      \begin{itemize}
        \item Temperatura: $-40$ a $85$\,°C ($\pm 0.5$\,°C)
        \item Humedad: 0 -- 100\,\%\,HR ($\pm 3$\,\%)
        \item Presión: 300 -- 1100\,hPa ($\pm 1$\,hPa)
      \end{itemize}
    \end{column}
    \begin{column}{0.44\textwidth}
      \begin{center}
        \fcolorbox{GrisLinea}{GrisClaro}{%
          \begin{minipage}[c][5.2cm][c]{0.9\textwidth}
            \centering
            {\color{GrisMedio}\small Diagrama de planta\\[0.4em]
            [insertar imagen aquí]\\[0.4em]
            \FuenteCodigo\tiny diagrama\_nodos.png}
          \end{minipage}%
        }
      \end{center}
    \end{column}
  \end{columns}
\end{frame}

%% -----------------------------------------------------------------
%% DIAPOSITIVA 6: CÓDIGO  [fragile — requerido por lstlisting]
%% -----------------------------------------------------------------
\begin{frame}[fragile]{Firmware ESP32 --- Publicación MQTT}
  \begin{tcolorbox}[cajaCode, top=2pt, bottom=2pt]
\begin{lstlisting}[language=Python]
from umqtt.simple import MQTTClient
from machine import I2C, Pin
import bme280, json, time

BROKER = "192.168.1.100"
TOPIC  = b"coriot/nodo1/ambiental"

def leer_y_publicar(client, sensor):
    t, p, h = sensor.read_compensated_data()
    datos = {"temp": t/100, "hum": h/1024, "pres": p/25600}
    client.publish(TOPIC, json.dumps(datos).encode(), qos=1)

i2c    = I2C(0, sda=Pin(21), scl=Pin(22), freq=400000)
sensor = bme280.BME280(i2c=i2c)
client = MQTTClient("nodo1", BROKER)
client.connect()
while True:
    try:
        leer_y_publicar(client, sensor)
    except OSError:
        client.connect()
    time.sleep(5)
\end{lstlisting}
  \end{tcolorbox}
  {\small\color{GrisMedio}\itshape
    Script MicroPython en cada nodo. QoS\,1 garantiza entrega al broker MQTT.}
\end{frame}

%% -----------------------------------------------------------------
%% DIAPOSITIVA 7: RESULTADOS
%% -----------------------------------------------------------------
\begin{frame}{Resultados de Desempeño}
  \begin{columns}[T]
    \begin{column}{0.32\textwidth}
      \begin{tcolorbox}[cajaIndicador]
        {\FuenteTitulo\LARGE\color{AzulITM}\textbf{98.7\%}}\\[4pt]
        {\small\color{GrisMedio}Disponibilidad}
      \end{tcolorbox}
    \end{column}
    \begin{column}{0.32\textwidth}
      \begin{tcolorbox}[cajaIndicador]
        {\FuenteTitulo\LARGE\color{AzulITM}\textbf{$\pm$0.5\,°C}}\\[4pt]
        {\small\color{GrisMedio}Precisión}
      \end{tcolorbox}
    \end{column}
    \begin{column}{0.32\textwidth}
      \begin{tcolorbox}[cajaIndicador]
        {\FuenteTitulo\LARGE\color{AzulITM}\textbf{5.2\,s}}\\[4pt]
        {\small\color{GrisMedio}Latencia promedio}
      \end{tcolorbox}
    \end{column}
  \end{columns}

  \vspace{0.6em}
  {\small
  \rowcolors{2}{GrisClaro}{white}
  \begin{tabular}{p{3.8cm}
      >{\centering\arraybackslash}p{1.7cm}
      >{\centering\arraybackslash}p{1.7cm}
      >{\centering\arraybackslash}p{1.7cm}
      p{3.2cm}}
    \rowcolor{AzulITM}
    \color{Blanco}\FuenteTitulo\scriptsize\bfseries Indicador &
    \color{Blanco}\FuenteTitulo\scriptsize\bfseries Meta &
    \color{Blanco}\FuenteTitulo\scriptsize\bfseries Logrado &
    \color{Blanco}\FuenteTitulo\scriptsize\bfseries \% Cumpl. &
    \color{Blanco}\FuenteTitulo\scriptsize\bfseries Observación \\
    Disponibilidad       & $\geq$97\%      & 98.7\%      & 101.8\% & Supera meta. \\
    Precisión temperatura & $\pm$1\,°C     & $\pm$0.5\,°C & 100\%  & Validada con patrón. \\
    Nodos activos        & 8               & 8            & 100\%   & Todos operativos. \\
  \end{tabular}
  }
\end{frame}

%% -----------------------------------------------------------------
%% DIAPOSITIVA 8: CIERRE  [plain → sin footline]
%% -----------------------------------------------------------------
\begin{frame}[plain]
  \begin{tikzpicture}[remember picture, overlay]
    \fill[AzulITM] (current page.south west) rectangle (current page.north east);
    \fill[AzulDigital]
      (current page.south west) rectangle ([yshift=0.5cm] current page.south east);
    \fill[AzulDigital]
      ([xshift=-2.4cm, yshift=-1.0cm] current page.north east)
      rectangle ([xshift=-1.9cm, yshift=-1.5cm] current page.north east);
    \fill[Magenta]
      ([xshift=-1.8cm, yshift=-1.5cm] current page.north east)
      rectangle ([xshift=-1.4cm, yshift=-1.9cm] current page.north east);
  \end{tikzpicture}

  \begin{center}
    \vspace{1cm}
    \LogoCoriotTikZ{AzulDigital}{white}\\[1.2em]
    {\FuenteTitulo\Huge\color{white}\textbf{¡Gracias!}}\\[0.4em]
    {\color{AzulDigital}\rule{0.4\textwidth}{1.5pt}}\\[0.5em]
    {\FuenteCodigo\normalsize\color{AzulDigital}%
      CONTROL\ \ \textbullet\ \ ROB\'OTICA\ \ \textbullet\ \ IOT}\\[1.2em]
    {\small\color{white}www.coriotlab.co\ \ |\ \ info@coriotlab.co}
  \end{center}
\end{frame}

\end{document}
